\section{Inflation Schedule Details}
\label{sec:inflation}

\subsection{Epoch-Level Emission Table}

The following table provides a detailed view of the emission schedule over the first 10 years:

\begin{table}[H]
\centering
\caption{Detailed emission schedule by year.}
\label{tab:emissions}
\begin{tabular}{rrrrr}
\toprule
\textbf{Year} & \textbf{$R_t$ (M LUX/yr)} & \textbf{Vesting (M LUX)} & \textbf{Est. Burn (M LUX)} & \textbf{Net (M LUX)} \\
\midrule
1 & 5.00 & 43.75 & 2.50 & +46.25 \\
2 & 4.85 & 43.75 & 5.00 & +43.60 \\
3 & 4.70 & 43.75 & 10.00 & +38.45 \\
4 & 4.56 & 43.75 & 15.00 & +33.31 \\
5 & 4.43 & 0 & 20.00 & $-15.57$ \\
6 & 4.29 & 0 & 25.00 & $-20.71$ \\
7 & 4.16 & 0 & 30.00 & $-25.84$ \\
8 & 4.04 & 0 & 35.00 & $-30.96$ \\
9 & 3.92 & 0 & 40.00 & $-36.08$ \\
10 & 3.80 & 0 & 45.00 & $-41.20$ \\
\bottomrule
\end{tabular}
\end{table}

Note the inflection point at year 5 when vesting completes: the network transitions from net inflationary to net deflationary, assuming moderate transaction growth (15\% annually).

\subsection{Reward Half-Life}

The reward emission has a half-life of:
\begin{equation}
    t_{1/2} = \frac{\ln 2}{\ln(1/\alpha)} = \frac{0.693}{0.0305} \approx 22.7 \text{ years}
\end{equation}

After 23 years, rewards are approximately 50\% of initial; after 46 years, approximately 25\%. This slow decay ensures continued validator incentivization for decades while maintaining long-term deflationary trajectory.

\subsection{Treasury Management}

The Foundation Reserve and Ecosystem Fund are managed under strict governance rules:

\begin{algorithm}[H]
\caption{Treasury Disbursement Protocol}
\label{alg:treasury}
\begin{algorithmic}[1]
\Require Disbursement request $(\mathsf{amount}, \mathsf{purpose}, \mathsf{recipient})$
\State \textbf{require} DAO approval with $\geq 66\%$ vote
\State \textbf{require} $\mathsf{amount} \leq 0.5\%$ of remaining treasury per quarter
\State \textbf{require} 7-day timelock after approval
\State Disburse from appropriate allocation
\State Publish transaction details publicly
\end{algorithmic}
\end{algorithm}

The quarterly cap of 0.5\% ensures treasury longevity: even with maximum disbursement every quarter, the treasury lasts at least 50 years.

\subsection{Deflationary Acceleration Scenario}

Under a high-growth scenario (25\% annual transaction growth), the deflationary acceleration is significant:

\begin{proposition}[Aggressive Deflation Bound]
If transaction volume grows at $g = 25\%$ annually and the average gas price remains constant, the circulating supply reaches $S_{\text{floor}}$ at approximately epoch 350 (year 7). At this point, fee burning ceases and all fees are redistributed to validators, creating a new equilibrium where:
\begin{equation}
    \text{Validator Revenue} = R_t + F_t^{\text{total}} \quad (\text{all fees, no burn})
\end{equation}
This increases validator yields and may attract additional staking, further tightening circulating supply.
\end{proposition}

\subsection{Cross-Chain Token Flow Accounting}

When LUX is bridged to external chains (Ethereum, Cosmos), the bridge-locked tokens affect circulating supply calculations:

\begin{definition}[Effective Circulating Supply]
\begin{equation}
    S_t^{\text{eff}} = S_t - \text{Staked}_t - \text{BridgeLocked}_t - \text{TreasuryLocked}_t
\end{equation}
The effective circulating supply excludes tokens that are economically inactive (staked, bridge-locked, or treasury-controlled).
\end{definition}

This metric provides a more accurate view of available market liquidity than raw circulating supply.

